\documentclass[11pt]{article}
\usepackage[margin=1in]{geometry}
\usepackage{times}
\usepackage{graphicx}
\usepackage{booktabs}
\usepackage{hyperref}

\title{Replication and Draw-Aware Extension Plan for Football Outcome Prediction}
\author{Team Member 1 \and Team Member 2 \and Team Member 3}
\date{}

\begin{document}
\maketitle

\begin{abstract}
We replicate a football outcome prediction study that evaluates Random Forest, XGBoost, and AdaBoost under imbalanced and SMOTE-balanced settings. We provide a reproducible implementation with configuration-driven experiments, run artifacts, and transparent gap analysis. On a full sample-data validation run, XGBoost performs best in both settings, while draw outcomes remain comparatively difficult in several variants. We also propose a draw-aware extension combining parity-based uncertainty features, cost-sensitive training, and calibration.
\end{abstract}

\section{Introduction}
This project performs replication and extension planning for a recent football match prediction paper. Our goal is to reproduce the core method, quantify alignment gaps, and provide a non-trivial future-work extension.

\section{Related Work}
Classical machine learning and boosting methods have been widely applied to sports forecasting, especially for structured tabular data. Prior work highlights persistent challenges from class imbalance and weak draw-class separability.

\section{Background}
The source paper uses Random Forest, XGBoost, and AdaBoost to predict \{home\_win, draw, away\_win\}, with preprocessing and class-balancing via SMOTE.

\section{Replication Setup}
We implemented:
\begin{itemize}
  \item Missing-value imputation (median/mode)
  \item Winsorization and min-max scaling
  \item One-hot encoding for categorical columns
  \item RandomForest-based feature selection
  \item Two settings: imbalanced and SMOTE
\end{itemize}

\section{Results}
Our sample-data validation run shows XGBoost as top performer in both settings. Detailed metrics and confusion matrices are stored in the repository under \texttt{results/}.

\section{Gap Analysis}
Differences from paper-reported numbers are primarily attributed to dataset mismatch (synthetic validation dataset), feature-space assumptions, and incomplete hyperparameter details in the source narrative.

\section{Extension and Future Work}
We propose a draw-aware extension:
\begin{itemize}
  \item Add parity-based uncertainty features
  \item Use cost-sensitive optimization for draw misclassification
  \item Apply probability calibration and a draw-specific threshold
\end{itemize}
Primary success target is a draw-F1 gain of at least 0.05 with minimal global-accuracy tradeoff.

\section{Conclusion}
The project now includes a reproducible replication pipeline, documented analysis, and a concrete extension hypothesis ready for implementation and final reporting.

\bibliographystyle{plain}
\bibliography{references}
\end{document}
